\section{Exercise 1}

Tên được đọc vào buffer 100 byte trong \texttt{.data} bằng syscall 8. \texttt{\$a0} là địa chỉ buffer, \texttt{\$a1} là số ký tự tối đa:

\begin{lstlisting}[language=MIPS,numbers=none]
  li $v0, 8    # syscall 8: đọc chuỗi từ bàn phím
  la $a0, name # $a0 = địa chỉ buffer nhận chuỗi
  li $a1, 100  # $a1 = số ký tự tối đa
  syscall
\end{lstlisting}

Syscall 8 giữ lại ký tự xuống dòng ở cuối chuỗi, nên nếu in thẳng thì dấu \texttt{!} sẽ bị xuống dòng. Vòng lặp dưới duyệt từng byte bằng \texttt{lb}, dừng ở ký tự kết thúc chuỗi, và thay ký tự xuống dòng (\texttt{0x0A}) bằng \texttt{0}:

\begin{lstlisting}[language=MIPS,numbers=none]
  la $t0, name                      # $t0 trỏ vào đầu buffer

find_newline_char:
  lb $t1, 0($t0)                    # nạp byte đang xét
  beqz $t1, print_msg               # gặp '\0' thì hết chuỗi
  beq $t1, 10, replace_newline_char # gặp '\n' thì thay thế
  addi $t0, $t0, 1                  # sang byte kế tiếp
  j find_newline_char

replace_newline_char:
  sb $zero, 0($t0)                  # ghi 0 đè lên ký tự xuống dòng
\end{lstlisting}

Cuối cùng in ba chuỗi liên tiếp bằng syscall 4: \texttt{"Hello, "}, tên và dấu \texttt{"!"}:

\begin{lstlisting}[language=MIPS,numbers=none]
print_msg:
  li $v0, 4         # syscall 4: in chuỗi
  la $a0, hello_msg # in "Hello, "
  syscall

  li $v0, 4         # syscall 4: in chuỗi
  la $a0, name      # in tên vừa nhập
  syscall

  li $v0, 4         # syscall 4: in chuỗi
  la $a0, excl_mark # in dấu "!"
  syscall
\end{lstlisting}

Nhập \texttt{John}:

\begin{lstlisting}[language={},numbers=none]
Hello, John!
\end{lstlisting}
